\chapter{知识清单　Concepts}\label{appendix:知识清单}

    \section{基本概念}\label{app:知识点-基本概念}

        \subsection{电解质与非电解质}

            \begin{itemize}
                \item \textcolor{red}{电解质}：在水溶液\textbf{或}熔融状态下能导电的\textbf{化合物}（酸、碱、盐、水）。
                \item \textcolor{red}{非电解质}：在水溶液\textbf{和}熔融状态下都不能导电的化合物（如 \cemh{CO2}、 \cemh{SO2}、 \cemh{NH3}、 \cemh{CCl4}、蔗糖、酒精）。
                \item \textcolor{blue}{必须是化合物本身电离，而非与水反应后产生离子；是物质的性质，与所处状态无关。}
                \item 溶液的导电能力 \(\propto\) 离子浓度（取决于溶液浓度与电离程度）。
            \end{itemize}

        \subsection{强电解质与弱电解质}

            \begin{itemize}
                \item 强电解质：在水溶液中\textbf{完全电离}（强酸、强碱、绝大多数盐），电离方程式用「\cemh{\xleq{}}」；溶液中只存在离子。
                \item 弱电解质：在水溶液中\textbf{部分电离}（弱酸、弱碱、水），电离方程式用「\cemh{<=>}」；离子、分子共存。
                \item 常见强酸：\cemh{H2SO4}、\cemh{HNO3}、\cemh{HCl}、\cemh{HBr}、\cemh{HI}、\cemh{HClO4}。
                \item 常见强碱：\cemh{KOH}、\cemh{NaOH}、\cemh{Ba(OH)2}、\cemh{Ca(OH)2}。
                \item 离子化合物都是强电解质；共价化合物中，强极性键对应强酸（强电解质）、较强极性键对应弱电解质、弱极性键对应非电解质。
            \end{itemize}

    \section{溶液中的三大平衡}\label{app:知识点-三大平衡}

        \subsection{水的电离与溶液的酸碱性}

            \cemh{H2O <=> H+ + OH^-}，水的电离极微弱。

            \begin{itemize}
                \item 水的离子积 \(K_\text{W} = \cemh{[H+]} \cdot \cemh{[OH^-]}\)：\textcolor{red}{只随温度变化，温度升高增大}（\SI{25}{\degreeCelsius} 时 \(\num{1E-14}\)）。
                \item 任何水溶液中，\cemh{[H+]} 与 \cemh{[OH^-]} 此消彼长。
                \item 中性：\cemh{[H+]} = \cemh{[OH^-]}；酸性：\cemh{[H+]} \(>\) \cemh{[OH^-]}；碱性：\cemh{[H+]} \(<\) \cemh{[OH^-]}。
                \item \(\cemh{pH} = -\lg{}\cemh{[H+]}\)，适用范围 \numrange{0}{14}。
                \item \textcolor{red}{酸碱抑制水的电离}（水电离的量看 \cemh{[H+]}、\cemh{[OH^-]} 中小的那个）；\textcolor{red}{盐类水解促进水的电离}（看大的那个）。
                \item 测量方法：酸碱指示剂、pH 试纸（不可润湿）、pH 计（测定前用蒸馏水洗净探头）。
            \end{itemize}

        \subsection{弱电解质的电离平衡}

            \(\cemh{HA <=> H+ + A-}\)，当 \(v_\text{电离} = v_\text{结合}\) 时达到电离平衡。\textcolor{red}{既然是平衡，所有平衡原理均适用。}

            \begin{itemize}
                \item \textcolor{red}{越热越电离}（升温促进电离）；\textcolor{red}{越稀越电离}（稀释促进电离）。
                \item \textcolor{red}{同离子效应}：加入相同离子，抑制电离；加入能消耗离子的物质，促进电离。
                \item 弱酸弱碱溶液每稀释 10 倍，pH 改变小于 1。
                \item 电离常数 \(K_i\) 只受温度影响，\textcolor{red}{温度升高 \(K_i\) 增大}。
                \item 多元弱酸分步电离，\(K_\text{a1} \gg K_\text{a2} \gg K_\text{a3}\)（一般相差 \numrange{1E4}{1E5}），\textcolor{red}{酸性主要由第一步电离决定}。
            \end{itemize}

        \subsection{盐类水解}

            \textcolor{red}{有弱就水解，无弱不水解，越弱越水解，谁强显谁性。}

            \begin{itemize}
                \item 盐电离出的离子与水电离出的 \cemh{H+} 或 \cemh{OH^-} 结合生成\textbf{弱电解质}的反应；可逆、微弱，是中和反应的逆反应（吸热）。
                \item \textcolor{red}{越热越水解}；\textcolor{red}{越稀越水解}；\textcolor{red}{水解促进水的电离}。
                \item 书写水解离子方程式：水写分子式、中间用可逆、后无沉气析（除双水解）；多元弱酸根分步水解。
                \item 强酸弱碱盐：正离子水解显酸性；弱酸强碱盐：负离子水解显碱性；弱酸弱碱盐：相互促进，酸碱性由相对强弱决定。
                \item 酸式盐：\cemh{NaHSO3}、\cemh{NaHC2O4} 显酸性（电离 \(>\) 水解）；\cemh{NaHCO3}、\cemh{NaHS} 显碱性（水解 \(>\) 电离）；\cemh{NaHSO4} 完全电离显酸性。
                \item \textcolor{red}{双水解（完全水解）}：正离子 \cemh{Al^{3+}}、\cemh{Fe^{3+}}；负离子 \cemh{S^{2-}}（\cemh{HS-}）、\cemh{CO3^{2-}}（\cemh{HCO3-}）、\cemh{SO3^{2-}}、\cemh{SiO3^{2-}}、\cemh{ClO-}。如泡沫灭火器：\cemh{Al^{3+} + 3HCO3- \xleq{} Al(OH)3 v + 3CO2 ^}。
            \end{itemize}

            应用：

            \begin{itemize}
                \item 净水：\cemh{Al^{3+} + 3H2O <=> Al(OH)3}（胶体）\cemh{+ 3H+}；\cemh{Fe^{3+}} 同理。
                \item 配制 \cemh{FeCl3} 溶液：将 \cemh{FeCl3} 溶于稀盐酸中（防止水解）。
                \item 纯碱洗涤油污加热效果更好：\cemh{Na2CO3 + H2O <=> NaHCO3 + NaOH} 吸热，升温促进水解。
                \item 草木灰不宜与铵态氮肥混用（\cemh{CO3^{2-}} 水解促进 \cemh{NH4+} 水解，氨逸失）。
                \item 氯化铵溶液除铁锈：\cemh{NH4+ + H2O <=> NH3 . H2O + H+}。
                \item 蒸干 \cemh{AlCl3} 溶液得不到氯化铝（水解生成 \cemh{HCl} 逸出），灼烧得 \cemh{Al2O3}；蒸干 \cemh{Al2(SO4)3} 得硫酸铝晶体。
            \end{itemize}

        \subsection{沉淀溶解平衡}

            难溶电解质饱和溶液中 \(v_\text{溶解} = v_\text{沉淀}\)，是动态平衡。

            \begin{itemize}
                \item 溶度积 \(K_\text{sp}\)：\textcolor{red}{只受温度影响}（大多随温度升高而增大），不受离子浓度影响。
                \item \textcolor{red}{溶度积规则}：\(Q < K_\text{sp}\) 无沉淀析出；\(Q = K_\text{sp}\) 达平衡；\(Q > K_\text{sp}\) 有沉淀析出。
                \item \textcolor{red}{同离子效应}：加入含相同离子的易溶强电解质，沉淀溶解平衡逆向移动，溶解度减小，\textcolor{red}{沉淀剂常加过量}。
                \item 沉淀转化：一般溶解度小的易转化成溶解度更小的，\(K_\text{sp}\) 相差越大转化越完全；\(K_\text{sp}\) 接近时取决于离子浓度。
                \item 沉淀溶解的方法：生成弱电解质（水）、生成气体、生成配合物（如 \cemh{AgCl} 溶于氨水）、氧化还原。
            \end{itemize}

    \section{溶液中的微粒浓度比较}\label{app:知识点-微粒浓度}

        \subsection{三大守恒}

            \begin{itemize}
                \item \textcolor{red}{电荷守恒}：溶液中阳离子所带正电荷总数 \(=\) 阴离子所带负电荷总数（记得乘以电荷数）。
                \item \textcolor{red}{物料守恒}：某元素各种存在形态的浓度之和等于该元素的初始浓度（微粒要找全）。
                \item \textcolor{red}{质子守恒}：由电荷守恒与物料守恒叠加得到。
            \end{itemize}

        \subsection{比较的一般思路}

            \begin{enumerate}
                \item 若发生反应，先算混合后各溶质浓度，再拆成微粒浓度（暂不考虑电离和水解）。
                \item 考虑所有平衡的影响：先主过程、后次过程（一般：电离 \(>\) 水解；弱电解质电离、盐类水解都微弱，水的电离更弱）。
                \item 用守恒关系列式验证。
                \item 全是一价离子时，由电荷守恒结合溶液酸碱性可推离子浓度大小关系。
            \end{enumerate}

    \section{酸碱滴定}\label{app:知识点-滴定}

        \begin{itemize}
            \item 原理：\(c_\text{酸} V_\text{酸} N_\text{酸} = c_\text{碱} V_\text{碱} N_\text{碱}\)（\(N\) 为酸、碱元数）。
            \item 滴定管：零刻度在上，最小刻度 \SI{0.1}{\milli\litre}，估读至 \SI{0.01}{\milli\litre}（数据记录到小数点后两位）；洗涤分自来水、蒸馏水、所盛溶液润洗三步。
            \item 锥形瓶：只能自来水洗、蒸馏水润洗，\textcolor{red}{不能用待测液润洗}。
            \item 终点判断：滴入最后半滴，指示剂变色，\textcolor{red}{并保持半分钟不复原}。
            \item 指示剂选择：变色范围要在滴定突跃范围内；石蕊不宜选用；\textcolor{red}{酸滴碱用甲基橙，碱滴酸用酚酞}；强碱滴弱酸用酚酞，强酸滴弱碱用甲基橙（或甲基红）。
            \item 误差分析：导致标准液用量偏大的操作使测定结果偏高（如滴定管未润洗、锥形瓶用待测液润洗等）。
            \item 氧化还原滴定：高锰酸钾法需酸化（用 \cemh{H2SO4}，不用 \cemh{HCl}、\cemh{HNO3}），一般不需另加指示剂；碘量法用淀粉溶液作指示剂。
        \end{itemize}

    \section{物质结构与性质}\label{app:知识点-结构性质}

        \subsection{VSEPR 理论与杂化轨道}

            \begin{itemize}
                \item 价层电子对数：
                \[
                \begin{aligned}
                    价层电子对数 &= \text{成键}\upsigma\text{电子对数} + \text{孤电子对数}\\
                                 &= n + \frac{\text{A 的价电子数} - \text{离子电荷} - n \times \text{B 的单电子数}}{2}
                \end{aligned}
                \]
                \item VSEPR 模型：2 对——直线形；3 对——平面正三角形；4 对——正四面体形；5 对——三角双锥；6 对——正八面体。
                \item 无孤电子对时，分子构型即 VSEPR 模型；有孤电子对时，扣除孤电子对后为实际构型。
                \item 杂化方式：2 对——\cemh{sp}；3 对——\cemh{sp^2}；4 对——\cemh{sp^3}；\cemh{PF5}——\cemh{sp^3d}；\cemh{SF6}——\cemh{sp^3d^2}。
                \item 典型实例：直线形 \cemh{BeCl2}、\cemh{CO2}、\cemh{HCN}；平面正三角形 \cemh{BF3}、\cemh{SO3}、\cemh{CO3^{2-}}；正四面体形 \cemh{CH4}、\cemh{NH4+}；三角锥形 \cemh{NH3}、\cemh{NF3}；角形 \cemh{H2O}、\cemh{SO2}。
                \item 键角：\textcolor{red}{孤电子对斥力 \(>\) 成键电子对斥力}（lp-lp \(>\) lp-bp \(>\) bp-bp），故键角 \cemh{H2O}（\ang{104.5}）\(<\) \cemh{NH3}（\ang{107;18}）\(<\) \cemh{CH4}（\ang{109;28}）。
                \item 等电子体：原子数相同、价电子总数相同的分子或离子，具有相似的化学键类型和空间结构。
            \end{itemize}

        \subsection{分子极性与分子间作用力}

            \begin{itemize}
                \item 极性分子：正、负电荷中心不重合；非极性分子：重合。
                \item 判断归纳：仅由非极性键形成的分子一定非极性；由极性键形成的双原子分子一定极性；由极性键形成的多原子分子看有无对称中心；\cemh{AB_n} 型分子若 A 的价层电子都已成键，则为非极性分子。
                \item 常见非极性分子：\cemh{CO2}、\cemh{CS2}、\cemh{BeCl2}、\cemh{BF3}、\cemh{SO3}、\cemh{CH4}、\cemh{CCl4}、\cemh{C2H4}、\cemh{C2H2}、苯。
                \item 常见极性分子：\cemh{H2O}、\cemh{H2S}、\cemh{SO2}、\cemh{NH3}、\cemh{PCl3}、\cemh{HX}、\cemh{CH3Cl} 等。
                \item 相似相溶：极性分子易溶于极性溶剂，非极性分子易溶于非极性溶剂。
                \item 范德华力：分子间普遍存在，比化学键弱得多，无方向性和饱和性；组成结构相似时，相对分子质量越大越强。
                \item 氢键：\cemh{X-H\bond{...}Y}（\cemh{X}、\cemh{Y} 可相同可不同），\cemh{H} 与 \cemh{N}、\cemh{O}、\cemh{F} 相连才形成；本质是静电作用，有方向性和饱和性，比范德华力强、比化学键弱。
            \end{itemize}

        \subsection{晶体}

            \par\vspace{0.6em}
            \begin{minipage}[t]{0.47\textwidth}
                \textbf{四类晶体：}
                \begin{datas}[leftmargin=*]
                    \item 离子晶体：微粒——正、负离子；键——离子键（无方向性、饱和性）；性质——熔沸点高、硬度较大、熔融或溶于水导电；熔点影响因素——电荷 \(>\) 半径。
                    \item 共价晶体：微粒——原子；键——共价键（有方向性、无密堆积）；性质——熔沸点很高、硬度很大、一般不导电；熔点影响因素——键能；实例——金刚石、\cemh{Si}、\cemh{SiO2}、\cemh{SiC}、\cemh{BN}。
                \end{datas}
            \end{minipage}
            \hfill
            \begin{minipage}[t]{0.47\textwidth}
                \begin{datas}[leftmargin=*]
                    \item 分子晶体：微粒——分子；键——分子间作用力；性质——熔沸点低、硬度小、一般不导电；熔点影响因素——氢键、相对分子质量、分子极性。
                    \item 金属晶体：微粒——金属阳离子、自由电子；键——金属键；性质——金属光泽、延展性、导电、导热；熔点影响因素——电荷数、半径。
                \end{datas}
            \end{minipage}
            \par\vspace{0.6em}

            \textcolor{blue}{石墨是混合型晶体：层内 \cemh{sp^2} 杂化共价键 \(+\) 离域\(\uppi\)键，层间范德华力；质软、能导电、熔点高。}

            \textcolor{red}{晶胞均摊法：顶角 \(1/8\)、棱上 \(1/4\)、面心 \(1/2\)、体心 \(1\)。}

            \textbf{金属堆积方式：}
            \begin{datas}[leftmargin=*]
                \item 简单立方：1 个原子、配位数 6、空间利用率 \SI{52.36}{\percent}（\cemh{Po}）。
                \item 体心立方：2 个、配位数 8、\SI{68.02}{\percent}（\cemh{Na}、\cemh{K}、\cemh{Fe}）。
                \item 面心立方：4 个、配位数 12、\SI{74.05}{\percent}（\cemh{Cu}、\cemh{Ag}、\cemh{Au}）。
                \item 六方：6 个、配位数 12、\SI{74.05}{\percent}（\cemh{Mg}、\cemh{Zn}、\cemh{Ti}）。
            \end{datas}
            \par\vspace{0.6em}
            \begin{minipage}[t]{0.47\textwidth}
                \textbf{典型晶胞粒子数：}
                \begin{datas}[leftmargin=*]
                    \item \cemh{NaCl}：\cemh{Na+} 4、\cemh{Cl-} 4。
                    \item \cemh{CsCl}：\cemh{Cs+} 1、\cemh{Cl-} 1。
                    \item \cemh{CaF2}：\cemh{Ca^{2+}} 4、\cemh{F-} 8。
                    \item \cemh{ZnS}：\cemh{Zn^{2+}} 4、\cemh{S^{2-}} 4。
                \end{datas}
            \end{minipage}
            \hfill
            \begin{minipage}[t]{0.47\textwidth}
                \textbf{共价晶体结构：}
                \begin{datas}[leftmargin=*]
                    \item 金刚石：晶胞含 \cemh{C} 8 个，配位数 4，\cemh{sp^3} 杂化，键长 \SI{154}{\pm}、键角 \ang{109;28}；\SI{12}{\gram} 金刚石含 \cemh{C-C} 键 \(2N_\text{A}\) 根。
                    \item \cemh{SiO2}：晶胞含 \cemh{Si} 8、\cemh{O} 16；每个 \cemh{Si} 连 4 个 \cemh{O}；\SI{1}{\mole} \cemh{SiO2} 含 \cemh{Si-O} 键 \(4N_\text{A}\) 根。
                    \item 冰：\SI{1}{\mole} \cemh{H2O} 含氢键 \SI{2}{\mole}（氢键有方向性，冰的密度小于水）。
                \end{datas}
            \end{minipage}
            \par\vspace{0.6em}

    \section{化学反应与电能}\label{app:知识点-电化学}

        \subsection{氧化还原反应}

            \begin{itemize}
                \item \textcolor{red}{守恒律}：得电子总数 \(=\) 失电子总数（化合价升降守恒）。
                \item \textcolor{red}{强弱律}：氧化性——氧化剂 \(>\) 氧化产物；还原性——还原剂 \(>\) 还原产物。
                \item \textcolor{red}{转化律}：同种元素化合价只靠拢、不交叉（归中）。
                \item \textcolor{red}{先后律}：氧化性（还原性）更强的先反应。还原性：\cemh{I-} \(>\) \cemh{Fe^{2+}} \(>\) \cemh{Br-}；氧化性：\cemh{I2} \(<\) \cemh{Fe^{3+}} \(<\) \cemh{Br2}。
                \item 氧化性、还原性受\textcolor{red}{浓度}、\textcolor{red}{pH}、\textcolor{red}{温度}影响；\cemh{H+} 参与的反应，\cemh{c(H+)}\(\uparrow\) 氧化性\(\uparrow\)。
            \end{itemize}

        \subsection{配平}

            \textcolor{red}{五步法}：\Circled{1}标出变化元素化合价\Circled{2}标出化合价升降数目（下角标计入）\Circled{3}求最小公倍数定系数\Circled{4}观察法确定其他系数\Circled{5}用未使用的原子检查。

            \textcolor{blue}{半反应书写：确定两极反应物产物 \(\to\) 得失电子数相同 \(\to\) 按电性补 \cemh{OH^-}、\cemh{H+} \(\to\) 书写。}

            \textcolor{blue}{离子方程式配平顺序：电子守恒 \(\to\) 电荷守恒 \(\to\) 原子守恒。}

        \subsection{原电池}

            \begin{itemize}
                \item Zn-Cu 原电池：\cemh{Zn} 为负极（失电子被氧化），\cemh{Cu} 为正极（\cemh{H+} 得电子）。
                \item \textcolor{red}{正正负负}：正离子 \(\to\) 正极，负离子 \(\to\) 负极；电子由负极经导线流向正极。
                \item 构成条件：两个活泼性不同的导电电极（活泼的作负极）、电极间闭合回路、电解质溶液（或熔融电解质）、自发的氧化还原反应。
                \item 盐桥作用：离子导体，平衡半电池电荷，使电流稳定。
                \item 离子交换膜：正离子交换膜、负离子交换膜、质子交换膜（只允许 \cemh{H+}）。
                \item 浓差电池：利用电解质浓度差异产生电势。
            \end{itemize}

        \subsection{化学电源}

            \begin{itemize}
                \item 锌锰干电池：普通（\cemh{NH4Cl}、\cemh{MnO2} 吸收 \cemh{H2} 防止极化）；碱性（用 \cemh{KOH}，比能量更高）。
                \item 铅蓄电池（二次电池）：\cemh{Pb + PbO2 + 4H+ + 2SO4^{2-} \xleq{} 2PbSO4 + 2H2O}，可根据硫酸密度判断充电程度。
                \item 氢氧燃料电池：介质不同负极反应不同——
                \begin{itemize}
                    \item 酸性：负极 \cemh{2H2 - 4e- \xleq{} 4H+}；正极 \cemh{O2 + 4H+ + 4e- \xleq{} 2H2O}。
                    \item 碱性：负极 \cemh{2H2 + 4OH^- - 4e- \xleq{} 4H2O}；正极 \cemh{O2 + 2H2O + 4e- \xleq{} 4OH^-}。
                    \item 熔融氧化物：负极 \cemh{2H2 + 2O^{2-} - 4e- \xleq{} 2H2O}；正极 \cemh{O2 + 4e- \xleq{} 2O^{2-}}。
                    \item 碳酸盐：负极 \cemh{2H2 + 2CO3^{2-} - 4e- \xleq{} 2H2O + 2CO2}；正极 \cemh{O2 + 2CO2 + 4e- \xleq{} 2CO3^{2-}}。
                \end{itemize}
                \textcolor{red}{极区 pH 变化看电极反应，不看总反应。}
            \end{itemize}

        \subsection{电解池}

            \begin{itemize}
                \item 阴极——与外电源负极相连（阳离子放电）；阳极——与外电源正极相连（阴离子放电）。
                \item 离子放电顺序见附录~\ref{appendix:常用数据}；\textcolor{blue}{活性电极（除 \cemh{Pt}、\cemh{Au}）优先于阴离子放电——阳极溶解}。
                \item \textcolor{red}{电解使非自发的氧化还原反应得以发生。}
                \item 应用：精炼铜（阳极粗铜溶解、阴极纯铜沉积）、电镀（镀层金属作阳极，电镀液浓度不变）、氯碱工业（\textcolor{red}{阳离子交换膜防止 \cemh{Cl2} 与 \cemh{NaOH} 反应}）。
                \item \textcolor{blue}{极区 pH 变化看电极反应，不等于溶液 pH 变化。}
            \end{itemize}

        \subsection{金属腐蚀与防护}

            \begin{itemize}
                \item 电化学腐蚀更普遍，常见的是\textcolor{red}{吸氧腐蚀}：负极 \cemh{Fe - 2e- \xleq{} Fe^{2+}}，正极 \cemh{O2 + 2H2O + 4e- \xleq{} 4OH^-}，铁锈 \cemh{Fe2O3 . xH2O}。
                \item 析氢腐蚀（酸性环境）：正极 \cemh{2H+ + 2e- \xleq{} H2 ^}。
                \item 防护：\textcolor{red}{牺牲阳极的阴极保护法}（船绑锌）、外加电流的阴极保护法（与电源负极相连）。
            \end{itemize}
